%% kapa: exam=1 round=20 year=2009 period=2 form=A subject=accounting subject_ko=회계학 seq=33 number=33
%% source: Q-Net 감정평가사 기출문제 (cst00309) artlSeq=5024178 file=extracted/1008/5024178/2009년 제20회 감정평가사 1차시험 기출문제/2교시.hwp
%% conversion: kordoc 4.24.51 -> tools/gen1.py (auto); review=proofread-sonnet; flags=none
\begin{question}{33}
종합원가계산을 적용하고 있는 (주)한국의 당기 원가자료는 다음과 같다.

{\small\centering
\begin{tabular}{@{}l r r r r r@{}}
\hline
 & & \multicolumn{2}{c}{직접재료원가} & \multicolumn{2}{c}{가공원가} \\
\cline{3-4}\cline{5-6}
\multicolumn{1}{c}{구분} & \multicolumn{1}{c}{물량} & \multicolumn{1}{c}{금액} & \multicolumn{1}{c}{완성도} & \multicolumn{1}{c}{금액} & \multicolumn{1}{c}{완성도} \\
\cline{1-1}\cline{2-2}\cline{3-3}\cline{4-4}\cline{5-5}\cline{6-6}
기초재공품 & 100개 & ￦5,000 & 100\% & ￦8,200 & 40\% \\
당기생산착수 & 900개 & ￦46,000 & - & ￦70,000 & - \\
기말재공품 & 200개 & ? & 100\% & ? & 60\% \\
당기완성품 & 800개 & ? & & ? & \\
\hline
\end{tabular}\par}
원가흐름 가정이 평균법이라면 당기완성품 제조원가는 얼마인가? \par
\begin{choices}
  \item ￦103,200
  \item ￦103,360
  \item ￦105,400
  \item ￦108,800
  \item ￦110,200
\end{choices}
\end{question}
\answer{4}
